%
%
\chapter{Understanding Unconferences}
This chapter defines what an Unconference is, briefly covering its origins, and describing what happens during a typical Unconference. These insights directly influence UnconfRS. Understanding Unconferences helps create a tool to aid in organizing them.

\section{Defining Unconferences}
An Unconference is a gathering of like-minded individuals who organize on the spot. The content of the Unconference is driven by the collective interests and expertise of the participants\cite{budd2015, duke2022}. In contrast, the typical conference is planned well in advance by few people, with the speakers, topics, and timing already determined. When presentations happen, there is an expectation for people to remain seated, listen, and hold questions until the end.

Unconferences tend to draw upon Open Space Technology (OST), a specific format of Unconference developed by Harrison Owen in the 1980s after the conclusion of the conference he put on revealed that the most useful parts were the coffee breaks \cite{harrison2008}. Some may use "Unconference" and "OST" interchangeably, but these are not the same. OST is a specific way of organizing and facilitating, while Unconference is a broader term not tied to a specific format. OST can be understood as an influential and popular type of Unconference. An OST meeting starts with participants in a circle with a blank agenda. Then they perform the agenda-setting exercise, putting their ideas on the wall (community bulletin board). After that, people gather in small groups (marketplace) for the ideas they are most interested in.\cite{harrison2008}

Some notable Unconferences include Foo Camp, Bar Camp, Ed Camp, and GSoC (Google Summer of Code) Mentor Summit. Foo Camp is an invitation-only event focused on computer science that started in 2003\cite{oreilly2018}. Bar Camp came about as an open alternative to Foo Camp in 2005 after Tantek Celik hadn't (yet) received an invitation to Foo Camp\cite{celik2006}. Ed Camp adopted the Unconference model for educators for professional development\cite{edcamp2011}. GSoC Mentor Summit adopted the Unconference model for bringing together mentors following the conclusion of that year's GSoC program\cite{googleblog2006}.

\section{A Typical Unconference}
A typical Unconference has three phases: agenda creation, breakout sessions, and gathering together. Agenda creation involves participants proposing session topics, usually written on post-it notes or index cards placed on a flipboard or whiteboard. Once the agenda is created and scheduled, these breakout sessions discuss the proposed topics. The event concludes with reconvening to share what has been learned.

Breakout sessions are distributed across multiple rooms and run in short consecutive timeslots. Multiple sessions are run in parallel, with participants being able to move freely between them. The schedule is posted visibly, allowing attendees to plan which sessions to attend.

Results are captured through session notes taken during discussions. At the closing gathering, attendees summarize sessions and share what they've learned throughout the experience.